\chapter{Modular arithmetic}\label{ch:mod}
\index{modular arithmetic}

Chapter~5 ended on a remainder. This chapter makes remainders the whole
number system: when ``$a$ and $b$ leave the same remainder mod $m$'' becomes
an equality --- congruence --- arithmetic collapses to a finite clock where
enormous questions (the last digits of $7^{2026}$; the day of week of any
date; which congruences have solutions at all) have finite answers. After
this chapter you will be able to compute powers by cycle, solve linear
congruences by the gcd-gate you already own, and read every AMC ``last
digit'' item as a period-hunting problem.

\section{Congruence as an equivalence}\label{sec:cong}
\index{congruence}

\begin{definition}
For a modulus $m \ge 1$, write $a \equiv b \pmod m$ (``$a$ is congruent to
$b$ mod $m$'') when $m \mid (a - b)$ --- equivalently, when $a$ and $b$
leave the same remainder under the division algorithm.
\end{definition}

Both clauses matter: the $m \mid (a-b)$ form is what proofs unfold; the
same-remainder form is what intuition runs on. Congruence is an equivalence
relation on $\Z$ (reflexive: $m \mid 0$; symmetric: $m \mid (a-b)
\Rightarrow m \mid (b-a)$; transitive: $m$ divides a sum of multiples) ---
and Figure~\ref{fig:classes} draws what that bought us for $m = 4$: a
partition of all of $\Z$.

\begin{figure}[ht]\centering
% Figure: partition of Z mod 4 --- four residue classes as bars.
% Source for Figure~\ref{fig:classes}. TikZ only; regenerable by compiling main.tex.
\begin{tikzpicture}[
  bar/.style={draw=rulegray,minimum height=7mm,font=\scriptsize,align=left},
  lab/.style={font=\footnotesize\sffamily,text=slate},
]
\node[bar,anchor=west,fill=accentsoft,text width=11.5cm] at (0,0)
  {$[0]_4=\{\dots,-8,-4,\mathbf{0},4,8,12,16,\dots\}$};
\node[bar,anchor=west,text width=11.5cm] at (0,-0.8)
  {$[1]_4=\{\dots,-7,-3,\mathbf{1},5,9,13,17,\dots\}$};
\node[bar,anchor=west,fill=accentsoft,text width=11.5cm] at (0,-1.6)
  {$[2]_4=\{\dots,-6,-2,\mathbf{2},6,10,14,18,\dots\}$};
\node[bar,anchor=west,text width=11.5cm] at (0,-2.4)
  {$[3]_4=\{\dots,-5,-1,\mathbf{3},7,11,15,19,\dots\}$};
\node[lab] at (6.2,-3.2) {four classes, zero overlap, covering all of $\Z$ ---
 $a\equiv b\pmod4$ iff they sit in the same bar};
\end{tikzpicture}

\caption{$\Z$ partitioned mod $4$ into residue classes
$[r]_4 = \{n : n \equiv r\}$. Notice the shape of equivalence itself: four
classes, no overlap, nothing left out --- Chapter~4's machinery doing real
work.}
\label{fig:classes}
\end{figure}

\begin{theorem}[Arithmetic respects congruence]\label{thm:ops}
If $a \equiv a' \pmod m$ and $b \equiv b' \pmod m$, then
$a + b \equiv a' + b'$, $a - b \equiv a' - b'$, and $ab \equiv a'b'
\pmod m$.
\end{theorem}

The proof is one element-chase ($a - a' = mk$, expand) --- but the content
is the license: \emph{inside any $+$, $-$, $\times$ computation you may swap
any number for anything congruent to it.} Reduce first, compute second,
always.

\begin{remark}
Division is conspicuously absent from Theorem~\ref{thm:ops}. That is not an
oversight --- it is the subject of Section~\ref{sec:cancel}, and forgetting
it is this chapter's most expensive mistake.
\end{remark}

Figure~\ref{fig:clock} makes the arithmetic physical: mod $12$ is the clock
face you already know, and adding $5$ hours is an arrow around it.

\begin{figure}[ht]\centering
% Figure: residue clock mod 12, with +5 and +7 step arrows.
% Source for Figure~\ref{fig:clock}. TikZ only; regenerable by compiling main.tex.
\begin{tikzpicture}[
  nd/.style={circle,draw,minimum size=6mm,font=\footnotesize},
  stp/.style={->,thick,accent},
  lab/.style={font=\footnotesize\sffamily,text=slate},
]
\foreach \i in {0,...,11}
{
  \node[nd] (n\i) at ({90 - \i*30}:2.1) {\i};
}
% a +5 walk from 2: 2->7->0->5->10, then +7: 10->5
\draw[stp,bend left=18] (n2) to (n7);
\draw[stp,bend left=18] (n7) to (n0);
\draw[stp,bend left=18] (n0) to (n5);
\draw[stp,bend left=18] (n5) to (n10);
\draw[stp,dashed,bend right=18] (n10) to (n5);
\node[lab] at (0,-3.0) {$+5$ steps: $2\to7\to0\to5\to10$ (solid); then $+7$ (dashed): $10\to5$};
\node[lab] at (0,-3.4) {arithmetic on the circle is arithmetic on residues};
\end{tikzpicture}

\caption{The residue clock mod $12$. Notice the wrap: $7 + 5$ lands on $0$,
then continues to $5$ --- and a $+7$ step from $10$ lands back on $5$ again.
Two different step sizes reaching the same class is the whole idea of
``solutions mod $m$''.}
\label{fig:clock}
\end{figure}

\begin{exercise}[easy]\label{exr:congdef}
Prove transitivity from the definition: $a \equiv b$ and $b \equiv c
\pmod m$ imply $a \equiv c \pmod m$.
\end{exercise}

\section{The residue system $\Z/m\Z$}\label{sec:zm}
\index{ZmZ@$\Z/m\Z$}

\begin{definition}
The \emph{least residue} of $a$ mod $m$ is the unique $r \in
\{0,1,\dots,m-1\}$ with $a \equiv r \pmod m$ --- the division-algorithm
remainder. The \emph{residue class} $[r]_m = \{n : n \equiv r \pmod m\}$,
and $\Z/m\Z$ is the set of all $m$ classes with the induced arithmetic.
\end{definition}

Computing ``mod $m$'' then means: do integer arithmetic, drop to least
residue at the end --- or, better, reduce \emph{as you go}, since
Theorem~\ref{thm:ops} says the answer never changes. $7^{100}$'s last digit
never requires the 85-digit number.

\begin{exercise}[easy]\label{exr:reduce}
Compute the least residue of $2^{10} + 3^{10}$ mod $5$ without evaluating
either power fully.
\end{exercise}

\section{Cancellation and multiplicative inverses}\label{sec:cancel}
\index{inverse!multiplicative}

Here is where $\Z/m\Z$ stops being a field in disguise and starts being a
ring worth studying. In $\Q$, $ab = ac$ implies $b = c$ for $a \ne 0$. Mod
$m$, it does not: $2 \cdot 3 \equiv 2 \cdot 0 \pmod 6$, yet $3 \not\equiv
0$. Figure~\ref{fig:units} shows why --- the multiplication table is a
permutation exactly when the row's factor is coprime to the modulus.

\begin{theorem}\label{thm:cancel}
$c$ has a multiplicative inverse mod $m$ --- some $x$ with $cx \equiv 1
\pmod m$ --- iff $\gcd(c, m) = 1$. Then $ca \equiv cb \pmod m$ implies
$a \equiv b \pmod m$.
\end{theorem}

\begin{proof}
Existence is B\'ezout, verbatim: $\gcd(c,m)=1 \Rightarrow cx + my = 1$, and
reading mod $m$ gives $cx \equiv 1$. Failure direction: if $d = \gcd(c,m)
> 1$, then $cx \equiv 1$ would force $d \mid 1$. For cancellation,
$ca \equiv cb$ means $m \mid c(a-b)$; with $\gcd(m,c)=1$, Euclid's lemma
gives $m \mid (a-b)$.
\end{proof}

\begin{figure}[ht]\centering
% Figure: multiplication tables mod 5 vs mod 6 --- units highlighted.
% Source for Figure~\ref{fig:units}. TikZ only; regenerable by compiling main.tex.
\begin{tikzpicture}[
  cell/.style={draw=rulegray,minimum width=6mm,minimum height=5mm,font=\scriptsize,inner sep=0.6pt},
  un/.style={cell,fill=accentsoft},
  hd/.style={cell,font=\scriptsize\bfseries},
  lab/.style={font=\footnotesize\sffamily,text=slate},
]
% mod 5 table rows 1..4
\node[lab,anchor=east] at (-0.7,0) {$\bmod\ 5$:};
\foreach \c/\h in {0/{},1/1,2/2,3/3,4/4}
  \node[hd] at (\c*0.62,0) {\h};
\foreach \r in {1,2,3,4}
{
  \node[hd] at (0,-\r*0.55) {\r};
  \foreach \c in {1,2,3,4}
  {
    \pgfmathtruncatemacro{\v}{mod(\r*\c,5)}
    \node[un] at (\c*0.62,-\r*0.55) {\v};
  }
}
\node[lab,anchor=west] at (3.3,-1.1) {every nonzero row is a permutation:};
\node[lab,anchor=west] at (3.3,-1.5) {each has a 1-entry $\Rightarrow$ inverse};
% mod 6 table rows 1..5
\node[lab,anchor=east] at (-0.7,-3.0) {$\bmod\ 6$:};
\foreach \c/\h in {0/{},1/1,2/2,3/3,4/4,5/5}
  \node[hd] at (\c*0.62,-3.0) {\h};
\foreach \r in {1,...,5}
{
  \node[hd] at (0,-3.0-\r*0.55) {\r};
  \foreach \c in {1,...,5}
  {
    \pgfmathtruncatemacro{\v}{mod(\r*\c,6)}
    \pgfmathtruncatemacro{\g}{gcd(\r,6)}
    \ifnum\g=1
      \node[un] at (\c*0.62,-3.0-\r*0.55) {\v};
    \else
      \node[cell] at (\c*0.62,-3.0-\r*0.55) {\v};
    \fi
  }
}
\node[lab,anchor=west] at (4.4,-3.8) {shaded rows ($\gcd=1$): still};
\node[lab,anchor=west] at (4.4,-4.2) {permutations $\Rightarrow$ inverses.};
\node[lab,anchor=west] at (4.4,-4.8) {unshaded rows collapse:};
\node[lab,anchor=west] at (4.4,-5.2) {no 1-entry, and $0$ appears ---};
\node[lab,anchor=west] at (4.4,-5.6) {$2\cdot3\equiv0$ kills cancellation.};
\end{tikzpicture}

\caption{Multiplication mod $5$ vs mod $6$. Notice what coprimality buys:
mod $5$ every nonzero row is a permutation (a 1 in every row: inverses for
all). Mod $6$, only the $\gcd = 1$ rows permute; the row of $2$ repeats
$\{0,2,4\}$ with no 1 --- and the $0$-entries are zero divisors, the reason
cancellation fails.}
\label{fig:units}
\end{figure}

Inverses are computed, not guessed: run Euclid backwards
(Example~\ref{ex:bezout}) and take the coefficient of $c$.

\begin{exercise}[medium]\label{exr:inverse}
Find $7^{-1} \pmod{10}$ and $8^{-1} \pmod{13}$ by extended Euclid; verify
each by multiplying back.
\end{exercise}

\section{Linear congruences}\label{sec:lincong}
\index{linear congruence}

\begin{theorem}
$ax \equiv b \pmod m$ has solutions iff $\gcd(a,m) \mid b$; then it has
exactly $\gcd(a,m)$ solutions mod $m$: divide through by $d =
\gcd(a,m)$ to reach $\frac{a}{d}x \equiv \frac{b}{d} \pmod{\frac{m}{d}}$
(coprime now), invert, and lift the answer back by $+k\cdot\frac{m}{d}$.
\end{theorem}

This \emph{is} Theorem~\ref{thm:dioph} in disguise --- $ax \equiv b
\pmod m$ literally means $ax - my = b$.

\begin{example}\label{ex:lincong}
Solve $14x \equiv 6 \pmod{20}$.
\par\smallskip
\textbf{Strategy.} Gcd first --- never cancel before checking.
$\gcd(14,20)=2$ and $2 \mid 6$, so solutions exist and there are exactly
two mod $20$. Divide the \emph{whole} congruence (modulus included) by $2$,
then invert the now-coprime coefficient.
\par\smallskip
\begin{solution}
Divide: $7x \equiv 3 \pmod{10}$. Since $7 \cdot 3 = 21 \equiv 1
\pmod{10}$, $7^{-1} \equiv 3$, so $x \equiv 3 \cdot 3 = 9 \pmod{10}$.
Lifting back, the two solutions mod $20$ are $x \equiv 9$ and $x \equiv
19$. Check: $14 \cdot 9 = 126 \equiv 6$ and $14 \cdot 19 = 266 \equiv 6
\pmod{20}$.
\end{solution}
\end{example}

\section{Modular patterns: last digits, power cycles, casting out
nines}\label{sec:patterns}
\index{power cycle}\index{casting out nines}

Powers of $a$ mod $m$ can produce at most $m$ residues, so the sequence
$a^1, a^2, a^3, \dots \pmod m$ eventually repeats --- and once it returns to
$1$ it loops with some \emph{period}: exponents differing by the period give equal powers.
``Last digit(s) of $a^N$'' is never a big-number problem: it is ``find the
period, then reduce $N$ mod the period.''

\begin{example}\label{ex:powers}
Find the last two digits of $7^{2026}$.
\par\smallskip
\textbf{Strategy.} Hunt the period first: compute $7^k \bmod 100$ until it
returns to $1$. Then the exponent only matters mod the period --- $2026$
reduces, $7^{2026}$ never has to exist.
\par\smallskip
\begin{solution}
$7^1 \equiv 07$, $7^2 \equiv 49$, $7^3 \equiv 343 \equiv 43$,
$7^4 \equiv 301 \equiv 01 \pmod{100}$. Period $= 4$. Since
$2026 = 4\cdot506 + 2$, $7^{2026} \equiv 7^2 \equiv 49 \pmod{100}$: the
last two digits are \textbf{49}.
\end{solution}
\end{example}

\begin{figure}[ht]\centering
% Figure: power cycle of 7 mod 100 --- a 4-node directed cycle.
% Source for Figure~\ref{fig:cycle}. TikZ only; regenerable by compiling main.tex.
\begin{tikzpicture}[
  nd/.style={circle,draw=accent,thick,minimum size=9mm,font=\small},
  ed/.style={->,thick,draw=rulegray},
  lab/.style={font=\footnotesize\sffamily,text=slate},
]
\node[nd] (a) at (0,0) {$01$};
\node[nd] (b) at (2.2,0) {$07$};
\node[nd] (c) at (2.2,-1.8) {$49$};
\node[nd] (d) at (0,-1.8) {$43$};
\draw[ed] (a) -- node[lab,above,midway] {$\times 7$} (b);
\draw[ed] (b) -- node[lab,right,midway] {$\times 7$} (c);
\draw[ed] (c) -- node[lab,below,midway] {$\times 7$} (d);
\draw[ed] (d) -- node[lab,left,midway] {$\times 7$} (a);
\node[lab] at (4.2,-0.9) {$7^1\equiv07$, $7^2\equiv49$,};
\node[lab] at (4.2,-1.3) {$7^3\equiv43$, $7^4\equiv01$};
\node[lab] at (4.2,-1.9) {period $=4$: exponent};
\node[lab] at (4.2,-2.3) {lives mod $4$};
\node[lab] at (1.1,-2.7) {$2026 \equiv 2 \pmod 4 \;\Rightarrow\; 7^{2026}\equiv 49$};
\end{tikzpicture}

\caption{The power cycle of $7$ mod $100$: four residues on a loop. Notice
where the exponent lives --- not on the circle but \emph{counting laps
around it}: $a^k$ depends only on $k \bmod 4$, which is why $2026 \equiv 2$
settles the question.}
\label{fig:cycle}
\end{figure}

\begin{checkpoint}[Predict before you run]
Listing~\ref{lst:cycle} finds that same cycle mechanically. Predict the
period for $(a,m) = (3,7)$ before reading the output below --- then check
what Fermat's little theorem (the olympiad-track version is in~\cite{zhao}) predicts.
\end{checkpoint}

\begin{lstlisting}[language=Python,caption={\texttt{powers\_cycle.py}:
multiply until a residue repeats (full script in
\texttt{code/powers\_cycle.py}).},label={lst:cycle}]
seen = {}; cur = 1 % m; k = 0
while cur not in seen:
    seen[cur] = k; cur = (cur * a) % m; k += 1
\end{lstlisting}
Real output (defaults $a=7$, $m=100$; then $3 \bmod 7$):
\begin{lstlisting}[numbers=none]
powers of 7 mod 100:
  7^0 == 1
  7^1 == 7
  7^2 == 49
  7^3 == 43
  7^4 == 1  (repeat)
cycle starts at k=0, period = 4
so 7^2026 mod 100 = 7^2 mod 100 = 49
---
powers of 3 mod 7:
  3^0 == 1
  3^1 == 3
  3^2 == 2
  3^3 == 6
  3^4 == 4
  3^5 == 5
  3^6 == 1  (repeat)
cycle starts at k=0, period = 6
so 3^2026 mod 7 = 3^4 mod 7 = 4
\end{lstlisting}

\begin{example}\label{ex:n3n}
Prove $n^3 - n$ is divisible by $6$ for every integer $n$.
\par\smallskip
\textbf{Strategy.} A residue bash is a \emph{finite} verification: there
are only six classes mod $6$, so check each once and the statement is
settled for all $n$ at once --- no inductive step to invent, no
factorisation to spot. The bash is the dumbest method in the chapter and
also the most reliable.
\par\smallskip
\begin{solution}
Read $n^3 - n \bmod 6$ for $n \equiv 0,1,2,3,4,5$:
$0 - 0 = 0$; $1 - 1 = 0$; $8 - 2 = 6 \equiv 0$; $27 - 3 = 24 \equiv 0$;
$64 - 4 = 60 \equiv 0$; $125 - 5 = 120 \equiv 0$.
Every class gives $0$, so $6 \mid n^3 - n$ for all $n$. (The structural
reason: $n^3 - n = (n-1)n(n+1)$, three consecutive integers --- one
divisible by $3$, at least one even.)
\end{solution}
\end{example}

Casting out nines is the same mechanism a level up: $10 \equiv 1 \pmod 9$
makes every integer congruent to its digit sum (Chapter~5 proved the
criterion); the checksum survives addition, so ``check the multiplication
by digit-summing both sides'' catches most slip errors. Note $9$'s blind
spot: it cannot see a transposition of digits that changes the answer by a
multiple of $9$ --- a checksum, not a proof.

\begin{exercise}[medium]\label{exr:cycle}
Find the last two digits of $3^{2026}$ (hint: the period mod $100$ is
longer than $4$ --- hunt it with \texttt{powers\_cycle.py}, then explain
why $\varphi(100) = 40$ is an upper bound).
\end{exercise}

\section{Two moduli at once: a CRT preview}\label{sec:crt}
\index{Chinese remainder theorem}

What if a number must satisfy \emph{two} congruences, $x \equiv 2 \pmod 3$
and $x \equiv 3 \pmod 5$? Hunting the finite list mod $15$ finds $x \equiv
8$ --- and the Chinese remainder theorem (Chapter~5's B\'ezout coefficients
are its engine; the olympiad track develops it in full~\cite{zhao}) says the
solution is unique mod $3\cdot5$ whenever the moduli are coprime. Chapter~5
solved one equation in integers; this is two conditions at once, and the
same gcd-gate decides existence.

\begin{example}\label{ex:calendar}
What day of the week was October 10, 2026?
\par\smallskip
\textbf{Strategy.} Count in mod $7$: January~1, 2026 was a Thursday (a
verifiable anchor). Count the day-of-year of October~10 --- a sum of month
offsets, not a loop through the calendar --- reduce mod $7$, and step that
many days from the anchor.
\par\smallskip
\begin{solution}
2026 is not a leap year. Days before October:
$31+28+31+30+31+30+31+31+30 = 273$, so October~10 is day $283$ of the year,
offset $282$ days from January~1. $282 = 7\cdot40 + 2$, so
$282 \equiv 2 \pmod 7$ and the day is Thursday $+ 2 =$
\textbf{Saturday}.
\end{solution}
\end{example}

\section{Chapter summary}
\begin{table}[ht]\centering
\begin{tabularx}{\linewidth}{@{}lYY@{}}
\toprule
\tablehead\thead{Object} & \thead{Rule} & \thead{Watch}\\
\midrule
$a \equiv b \pmod m$ & $m \mid (a-b)$; equivalence on $\Z$ & statement, not equation\\
$+, -, \times$ & substitute freely, reduce early & exponent is NOT reduced mod $m$\\
$ca \equiv cb$ & cancel iff $\gcd(c,m)=1$ & else divide modulus too\\
$c^{-1} \bmod m$ & B\'ezout coefficient of $c$ & exists iff coprime\\
$ax \equiv b$ & $\gcd \mid b$ gate; $\gcd$ solutions & divide modulus by $d$ as well\\
$a^N \bmod m$ & find period, reduce $N$ mod period & period $\le m$, usually $\varphi$\\
digit sum & $n \equiv$ digits mod $9$ & checksum, not proof\\
\bottomrule
\end{tabularx}
\caption{The chapter in one page. Every rule is a Chapter-5 theorem wearing
a clock face --- the gcd still decides everything.}
\label{tab:mod}
\end{table}

\section{Common mistakes}
\index{common mistakes}
\begin{itemize}
\item \textbf{Dividing by a non-unit.} $2x \equiv 4 \pmod 6$ ``solves'' to
$x \equiv 2$ --- but $x \equiv 5$ also works, and the real solution set is
$\{2,5\}$. Cancelling $2$ without dividing the modulus deleted half the
answer. Check $\gcd(c,m)$ \emph{before} cancelling.
\item \textbf{Writing $x = 3 \pmod 5$ inside algebra.} Congruence is a
relation between \emph{two} integers, not a decoration on one. The $x$
that equals $3$ is an integer; the $x$ congruent to $3 \bmod 5$ is a class.
Mixing them produces ``$3 = 8$'' nonsense that quietly poisons a proof.
\item \textbf{Reducing the exponent by the modulus.} $a^N \bmod m$ cycles
with the \emph{period} of $a$'s powers --- period $4$ for $7 \bmod{100}$,
not $100$. $7^{2026}$ is not $7^{26}$. Find the period first; only then
does the exponent get reduced.
\end{itemize}

\section{Additional exercises}
\begin{exercise}[easy]
Show $a \equiv b \pmod m$ iff $a$ and $b$ have the same least residue ---
and check whether $-7 \equiv 5 \pmod{12}$ (careful with the definition).
\end{exercise}
\begin{exercise}[easy]
Compute $17 \cdot 23 \pmod 5$, $9^{3} \pmod 7$, and
$2026 \cdot 10^{10} \pmod 9$ by early reduction.
\end{exercise}
\begin{exercise}[easy]
Which elements of $\Z/12\Z$ have inverses? List them with their inverses.
\end{exercise}
\begin{exercise}[easy]
Use casting out nines to catch the error in $47 \times 83 = 3911$ --- then
produce a wrong answer the check \emph{cannot} catch.
\end{exercise}
\begin{exercise}[medium]
Solve $12x \equiv 8 \pmod{28}$ completely, and prove $15x \equiv 7
\pmod{21}$ has no solutions --- state the gate that decides each.
\end{exercise}
\begin{exercise}[medium]
Prove $n^3 - n$ is divisible by $6$ for every integer $n$ by checking the
six residue classes --- then prove it again by factoring. Which proof does
the checker trust more, and why?
\end{exercise}
\begin{exercise}[medium]
Find the last digit of $2026^{2026}$ and the last two digits of
$9^{100}$ (period mod $4$ vs mod $100$).
\end{exercise}
\begin{exercise}[medium]
Prove the ``rule of $11$'' as a congruence ($10 \equiv -1$) and use it to
test $121\,121$ and $9\,876\,543$.
\end{exercise}
\begin{exercise}[hard]
Prove that the powers of $a \bmod m$ are \emph{purely} periodic exactly when
$\gcd(a,m)=1$; otherwise a tail precedes the cycle (find the tail for
$2 \bmod{12}$ and state the general rule).
\end{exercise}
\begin{exercise}[hard]
Solve $x \equiv 2 \pmod 3$, $x \equiv 3 \pmod 5$, $x \equiv 2 \pmod 7$
simultaneously. (Build it one congruence at a time --- this is the CRT
construction in embryo.)
\end{exercise}
\begin{exercise}[hard]
Extend \texttt{powers\_cycle.py} to print, for each base $a = 2,\dots,9$
and modulus $m$, the tail length and period --- and conjecture the rule
connecting them to $\gcd(a,m)$. Verify on $m = 12, 20, 100$.
\end{exercise}

\chapternotes
The equivalence-first presentation and the cancellation theorem follow
MIT 6.042J's modular-arithmetic sections~\cite{mcs}; Velleman's congruence
chapter supplied the residue-partition viewpoint~\cite{velleman}, and
Hammack's the clock-face intuition~\cite{hammack}. Berkeley CS70's
modular note is the recommended next read (it pushes straight into
Fermat--Euler, CRT and RSA~\cite{cs70}); Zhao's olympiad handout covers the
orders/primitive-roots continuation~\cite{zhao}. Contest items powered by
this chapter's machinery are on the AoPS archive~\cite{aops-amc}: power
cycles (AMC~12~2024~B~P14, AIME~2024~I~P13), last-digit and remainder
problems (AIME~2020~I~P10--12, AIME~2010~I~P2, AIME~2018~I~P11), linear
congruences (AIME~2023~I~P7, AIME~2025~II~P2) and CRT-flavoured counting
(AIME~2023~II~P15).
